\begin{textsample}{POS Dim 1 – 2009\_c6 – Score -2.00 – 2009\_c6/enade\_2009\_ci\_econ\_d1cc84d7e1.txt}  \label{ex:f1_pos_225}
Capítulo 3 Análise Técnica da Prova Este capítulo tem por objetivo apresentar o desempenho dos estudantes de CIÊNCIAS ECONÔMICAS no ENADE/2009.
Para isso, foram calculadas as estatísticas básicas da prova como um todo, bem como as estatísticas das questões relacionadas à Formação Geral e ao Componente Específico.
Nas Tabelas, são apresentadas as seguintes estatísticas : total da população, de presentes, média, erro-padrão da média, desvio-padrão, mediana, nota máxima, nota mínima e coeficiente de assimetria ( ver definição em 1.4.3 ).
Nas estatísticas das questões discursivas, foram detalhadas também as notas nulas : em branco, por protesto ou por insuficiência.
As estatísticas apresentadas neste capítulo contemplam, separadamente, os ingressantes, os concluintes e o total de estudantes, e foram calculadas tendo -se em vista as seguintes agregações : ( a ) Região Geográfica e País e ( b ) a Categoria Administrativa e a Organização Acadêmica.
As notas brutas das questões discursivas em cada uma das partes, Formação Geral/Componente Específico, foram calculadas a partir da média simples das notas de cada uma das questões correspondentes.
A primeira questão discursiva da Formação Geral foi desconsiderada e, portanto, a nota bruta referente às questões discursivas de Formação Geral foi igual à nota da questão mantida, a segunda.
No Componente Específico, a nota bruta correspondeu à média das notas nas três questões.
As notas brutas das questões objetivas de cada uma das partes, Formação Geral/Componente Específico, foram calculadas a partir dos acertos dos estudantes considerando -se uma regra de três para a conversão para a escala de 0 a 100, descontando -se aquelas questões que foram desconsideradas pela comissão ou que não passaram pelo teste do índice de discriminação ponto-bisserial ( ver seção 1.4.6 para a descrição do índice ).
As questões com índices de discriminação inferior a 0,2 foram excluídas do cômputo da nota do estudante.
As notas brutas de Formação Geral e do Componente Específico para cada um dos estudantes foram calculadas como médias ponderadas das suas notas nas questões objetivas e discursivas, ponderação diferente para cada parte : C\_FG = 0,6 × C\_FG\textasciicircum{}OBJ + 0,4 × C\_FG\textasciicircum{}DIS C\_CE = 0,85 × C\_CE\textasciicircum{}OBJ + 0,15 × C\_CE\textasciicircum{}DIS onde C\_FG é a nota bruta da Formação Geral, C\_FG\textasciicircum{}OBJ é a nota bruta nas questões objetivas da Formação Geral, C\_FG\textasciicircum{}DIS é a nota bruta nas questões discursivas da Formação Geral, C\_CE é a nota bruta do Componente Específico, C\_CE\textasciicircum{}OBJ é a nota bruta nas questões objetivas do Componente Específico, e C\_CE\textasciicircum{}DIS é a nota bruta nas questões discursivas do Componente Específico.
A nota bruta da prova, C\_GER, por sua vez, é também uma média ponderada da nota de cada uma das partes, Formação Geral/Componente Específico : C\_GER = C\_FG × 0,25 + C\_CE × 0,75. 3.1 Estatísticas Básicas da Prova 3.1.1 Estatísticas Básicas Gerais A Tabela 3.1 apresenta as estatísticas básicas da prova por grupo de estudantes.
O total de estudantes inscritos para a prova foi de 20821.
Destes, 75,3 \% compareceram, sendo que a presença foi maior entre concluintes ( 83,1 \% ) do que entre ingressantes ( 69,4 \% ).
A média geral da prova foi de 35,61, sendo que os ingressantes obtiveram média mais baixa, de 32,29, que os concluintes, cuja média foi 39,31.
O desvio-padrão geral foi de 14,9, sendo o do grupo de ingressantes menor ( 13,2 ), do que o do grupo dos concluintes ( 15,7 ), aqueles possuíram uma distribuição mais homogênea nas notas.
A nota máxima foi 89,9, obtida por pelo menos um concluinte ao passo que a maior nota obtida por um ingressante foi 81,1.
Os coeficientes de assimetria para os dois grupos considerados, o de concluintes e o de ingressantes, bem como para a população como um todo foram todos positivos, indicando uma maior concentração à esquerda da distribuição e um maior espalhamento à direita.
Uma comparação entre o desempenho de ingressantes e concluintes pode ser realizada a partir dos dados apresentados no Gráfico 3.1.
Cumpre destacar que as notas dos concluintes, de um modo geral, foram mais elevadas do que as dos ingressantes, cuja distribuição de notas está à esquerda dos primeiros.
A ideia de melhor desempenho é reforçada pelas notas medianas dos dois grupos : 38,5 entre os concluintes e 31,9 entre os ingressantes.
Nas faixas de notas até 40 pontos o percentual de ingressantes é maior que o dos concluintes, situação esta que se inverte para as faixas de notas superiores a 40 pontos.
Constatou -se que aproximadamente 10,1 \% dos ingressantes e 24,8 \% dos concluintes situaram -se nas faixas de notas superiores a 50,0 pontos.
O Gráfico 3.2 apresenta informações referentes ao desempenho geral de ingressantes e concluintes, comparando os resultados considerando a Região, Categoria Administrativa e Organização Acadêmica.
Levando -se em conta as notas médias dos estudantes em cada Região, observa -se que, nos dois grupos há uma variação expressiva entre as Regiões.
As notas médias mais elevadas entre os concluintes ocorreram nas Regiões Sudeste ( 40,86 ), Sul ( 40,25 ) e Centro-Oeste ( 39,32 ), todas acima da média nacional ( 39,30 ).
Em relação aos ingressantes, as médias mais altas foram obtidas nas Regiões Sudeste ( 34,38 ) nas demais Regiões a média foi inferior à média nacional ( 32,28 ).
Na Região Norte foi obtida a média de notas mais baixa entre os ingressantes ( 28,10 ) e os concluintes ( 32,74 ).
Cursos nesta área foram oferecidos em instituições federais, estaduais, municipais e privadas.
Levando -se em conta os agrupamentos dos estudantes em categorias administrativas, observa -se que a pontuação média mais elevada entre os concluintes e ingressantes foi encontrada em instituições de origem federal ( respectivamente 42,94 e 34,07 ) – ambas acima da média nacional.
Tendo como foco as notas médias dos estudantes agrupados por Organização Acadêmica, notamos que só universidades ofereceram as maiores médias para concluintes ( 40,24 ), ao passo que as Faculdades para ingressantes ( 33,22 ).
As piores notas médias foram observadas nos Centros Universitários, 35,72 para concluintes e 29,48 para ingressantes. 3.1.2 Estatísticas Básicas de Formação Geral A Tabela 3.2 apresenta as estatísticas básicas de Formação Geral por grupo de estudantes.
O total de estudantes inscritos para a prova foi de 20821.
Destes, 75,3 \% do total compareceram, sendo que a presença foi maior entre concluintes ( 83,1 \% ) do que entre ingressantes ( 69,4 \% ).
A média geral da parte da prova de Formação Geral foi de 47,89, sendo que os concluintes obtiveram média mais elevada, 49,01, que os ingressantes, cuja média foi 46,89.
O desvio-padrão geral foi de 21,0, sendo muito similar entre os grupos de ingressantes ( 21,1 ) e de concluintes ( 21,0 ).
A nota máxima foi 100, obtida por estudantes de ambos os grupos.
Para esta parte, assim como para a nota geral, os coeficientes de assimetria são positivos para todos os grupos.
Como, porém, os valores são bem perto de zero estes indicaram uma ligeira inclinação para a esquerda.
O Gráfico 3.3 permite a comparação entre concluintes e ingressantes relativamente ao desempenho na parte da prova que aborda a Formação Geral.
Nesta parte, se excluídas as provas com nota zero, tal como na prova como um todo, também os concluintes apresentaram desempenho superior, embora seja perceptível um equilíbrio.
A mediana dos concluintes foi maior ( 49,6 ) que a dos ingressantes ( 42,8 ), assim como a média.
Dos ingressantes, aproximadamente 45,2 \% alcançaram acima de 50 ; dentre os concluintes, o percentual dos estudantes nesta faixa de notas é 49,5 \%.
No Gráfico 3.4 são apresentadas as informações referentes ao desempenho dos ingressantes e dos concluintes na parte de Formação Geral, em diferentes agregações como Região, Categoria Administrativa e Organização Acadêmica.
Considerando -se as notas médias dos estudantes segundo cada Região, observou -se que, tanto no grupo dos ingressantes quanto no grupo dos concluintes, houve diferenças marcantes.
As médias variam, entre os concluintes, de 43,19 na Região Norte a 51,12 na Região Sul.
Entre os ingressantes, os valores foram de 42,08 na Região Norte a 49,96 na Região Sudeste.
A maior dispersão de notas dos estudantes de uma dada Região foi encontrada no Norte tanto para os concluintes ( 0,91 ) quanto para os ingressantes ( 1,01 ).
Na Região Sudeste, os concluintes ( 50,02 ) apresentaram médias quase iguais à dos ingressantes ( 49,96 ).
Em todas as Regiões, os concluintes apresentaram médias maiores do que os ingressantes.
Tendo como foco as categorias administrativas, as maiores notas médias foram das instituições estaduais para concluintes ( 52,55 ) e das instituições federais para ingressantes ( 48,34 ), enquanto às das instituições privadas foram menores para os concluintes ( 46,36 ) e para os ingressantes ( 45,40 ).
Os estudantes das instituições privadas se apresentam mais homogêneos do que as das demais instituições.
Considerando -se o tipo de Organização Acadêmica, Centros Universitários apresentaram os maiores valores para os concluintes ( 49,54 ) e Faculdades para os ingressantes ( 47,60 ).
Os Centros Federais de Educação Tecnológica e os Institutos de Educação, Ciência e Tecnologia não estavam representados nesta área. 3.1.3 Estatísticas Básicas do Componente Específico A Tabela 3.3 apresenta as estatísticas básicas referentes ao Componente Específico da prova, tendo sua média geral, de 31,57, inferior ao desempenho da prova de CIÊNCIAS ECONÔMICAS como um todo, cuja média foi 35,61, como apresentado na Tabela 3.1, que por sua vez foi inferior ao desempenho na Formação Geral, que teve média igual a 47,89, conforme mostrado na Tabela 3.2.
Os concluintes apresentaram um desempenho médio ( 36,12 ) maior ao dos ingressantes ( 27,47 ), o mesmo ocorrido para a nota da Formação Geral, havendo ainda indícios de maior homogeneidade entre as notas médias dos ingressantes do que entre as dos concluintes, como sugeriu o desvio padrão, menor no grupo dos ingressantes ( 13,8 ) do que no dos concluintes ( 17,1 ).
As notas máximas obtidas foram 82,6 para os ingressantes e 93,7 para os concluintes.
As distribuições de notas para cada grupo se apresentaram com uma inclinação para notas menores, como evidenciados pelos coeficientes de assimetria positivos.
A situação de assimetria é confirmada pelo Gráfico 3.5 que proporciona uma comparação do desempenho entre concluintes e ingressantes em relação ao Componente Específico da prova.
Claramente a distribuição das notas médias dos ingressantes, à direita daquela dos concluintes, evidenciou o percentual mais alto de notas médias nas faixas mais baixas para este grupo.
As medianas das distribuições das notas dos ingressantes e dos concluintes foram 26,5 e 35,1, respectivamente, enfatizando o melhor desempenho dos concluintes nesta parte da prova.
Assim como os Gráficos 3.2 e 3.4, o Gráfico 3.6 apresenta uma comparação dos resultados em relação às Regiões, Categoria Administrativa e Organização Acadêmica, agora levando em conta o desempenho dos ingressantes e dos concluintes no Componente Específico da prova.
Tendo como foco as notas médias dos estudantes, em determinada Região, observa -se que, entre os concluintes, as mais elevadas foram encontradas nas Regiões Sudeste ( 37,86 ) e Sul ( 36,68 ), valores acima da média nacional para esta parte ( 36,12 ).
Entre os ingressantes, as maiores notas ocorreram na Região Sudeste ( 29,24 ), acima da média nacional ( 27,46 ).
Novamente, para a análise segundo a Categoria Administrativa, todas as entidades foram representadas.
Observa -se que as médias mais elevadas, tanto entre os concluintes como entre os ingressantes, foram encontradas em instituições de origem federal ( 40,39 entre os concluintes e 29,36 entre os ingressantes ) enquanto nas instituições municipais as médias foram ( 32,03 entre os concluintes e 22,42 entre os ingressantes ).
Quanto à Organização Acadêmica, as Universidades, Centros Universitários e Faculdades são representadas com este curso.
Considerando -se o tipo de Organização Acadêmica, para os concluintes, Universidades apresentaram as maiores médias, seguidas das Faculdades e dos Centros Universitários.
Já para os ingressantes, as maiores médias se encontram nas Faculdades, seguidas das Universidades e Centros Universitários. 3.2 Análise das Questões Objetivas 3.2.1 Formação Geral A Tabela 3.4 apresenta as estatísticas básicas relativas às oito questões objetivas da Formação Geral dos estudantes.
Adiantamos que a primeira questão foi desconsiderada pelo critério da correlação ponto-bisserial.
Como pode ser observado na Tabela 3.4, os estudantes concluintes obtiveram desempenho maior, com nota média de 62,40, em relação aos estudantes ingressantes cuja nota média foi 60,79.
A variabilidade entre os concluintes foi um pouco menor do que a dos ingressantes, como indicaram os desvios-padrão de 22,5, para o grupo de concluintes e 23,0 para o grupo dos ingressantes.
A amplitude total das notas foi a mesma para ambos os grupos : nota mínima de 0 e máxima de 100.
A mediana foi igual para ambos os grupos : 57,1.
A assimetria, medida pelo coeficiente de assimetria ( skewness ) foi a mesma para os ingressantes e concluintes ( -0,5 ), ainda que ambos fossem negativos indicando uma maior concentração à direita.
A Tabela 3.5 apresenta a classificação das questões objetivas da Formação Geral, segundo o índice de facilidade ( questões em vermelho foram desconsideradas pelo critério ponto-bisserial ).
Como as questões foram comuns a todas as áreas, a tabela apresenta a classificação considerando -se os concluintes de todas as áreas e aqueles do curso de CIÊNCIAS ECONÔMICAS.
De acordo com os índices obtidos, as questões objetivas da prova foram assim avaliadas para a área de CIÊNCIAS ECONÔMICAS : das oito questões, três tiveram entre 60 \%, inclusive, e 85 \% de acertos e foram classificadas como fáceis ( somente uma para o total de concluintes ) ; A Tabela 3.6 descreve os resultados das análises das questões objetivas relativas à Formação Geral, segundo o poder de discriminação, utilizando -se para tal o índice de discriminação ponto-bisserial.
Como as questões foram comuns a todas as áreas, a tabela apresenta a classificação considerando -se os concluintes de todas as áreas e aqueles do curso de CIÊNCIAS ECONÔMICAS.
A desconsideração da questão segundo este critério considerou o universo de todos os concluintes.
Nenhuma das questões apresentou índices iguais ou acima de 0,40 ( nenhuma para a população de concluintes como um todo ), sendo, portanto, classificadas como muito boas nesta característica, para este grupo de estudantes, duas foram classificadas como boas com índice entre 0,30, inclusive, e 0,40 ( três para os concluintes como um todo ) e cinco foram classificadas como médias, com índice entre 0,20, inclusive, e 0,30 ( quatro para os concluintes como um todo ).
Uma das questões foi classificada como fraca tanto para a população de concluintes como um todo como para a de concluintes de CIÊNCIAS ECONÔMICAS e foi desconsiderada pelo critério ponto-bisserial. alcançando entre 40 \%, inclusive, e 60 \% de acertos ( três para o total de concluintes ) ; uma questão obteve entre 15 \%, inclusive, e 40 \% de acertos, sendo considerada difícil ( três para o total de concluintes ).
Por fim, uma foi classificada como muito difícil tento obtido percentual menor do que 15 \% de acertos ( uma para o total de concluintes ), coincidentemente a questão desconsiderada pelo critério ponto-bisserial.
Uma foi classificada na categoria muito fácil ( nenhuma para o total de concluintes ).
A Tabela 3.6 descreve os resultados das análises das questões objetivas relativas à Formação Geral, segundo o poder de discriminação, utilizando -se para tal o índice de discriminação ponto-bisserial.
Como as questões foram comuns a todas as áreas, a tabela apresenta a classificação considerando -se os concluintes de todas as áreas e aqueles do curso de CIÊNCIAS ECONÔMICAS.
A desconsideração da questão segundo este critério considerou o universo de todos os concluintes.
Nenhuma das questões apresentou índices iguais ou acima de 0,40 ( nenhuma para a população de concluintes como um todo ), sendo, portanto, classificadas como muito boas nesta característica, para este grupo de estudantes, duas foram classificadas como boas com índice entre 0,30, inclusive, e 0,40 ( três para os concluintes como um todo ) e cinco foram classificadas como médias, com índice entre 0,20, inclusive, e 0,30 ( quatro para os concluintes como um todo ).
Uma das questões foi classificada como fraca tanto para a população de concluintes como um todo como para a de concluintes de CIÊNCIAS ECONÔMICAS e foi desconsiderada pelo critério ponto-bisserial.
Na sequência, no Quadro 3.1, estão especificadas, por questão, as habilidades aferidas pelas questões de Formação Geral da prova de CIÊNCIAS ECONÔMICAS, bem como os resultados da análise relativa aos índices de facilidade e de discriminação de cada questão.
O índice de facilidade para a área de CIÊNCIAS ECONÔMICAS teve uma variação considerável, oscilando entre 0,11 até 0,89 e o de discriminação ficou entre 0,09 e 0,37.
As questões 6 e 7, que tiveram em comum aferir as habilidades de “ Analisar e criticar informações ”, “ Extrair conclusões por indução e/ou dedução ”, “ Estabelecer relações, comparações e contrastes em diferentes situações ”, “ Detectar contradições ”, “ Fazer escolhas valorativas avaliando consequências ”, “ Questionar a realidade “ e “ Argumentar coerentemente ” foram as que apresentaram, para a população de concluintes de CIÊNCIAS ECONÔMICAS, maior poder discriminatório ambas com índice de 0,37.
A questão mais fácil, com 89 \% de acertos, foi a de número 8, que aferiu também, além das já listadas anteriormente, a habilidade de “ Ler e interpretar textos ”.
De outra parte, a questão com menor índice de facilidade e de discriminação em Formação Geral foi a questão 1, que foi desconsiderada, e que deveria aferir as habilidades de “ Ler e interpretar textos ”, “ Analisar e criticar informações ”, “ Extrair conclusões por indução e/ou dedução ”, “ Estabelecer relações, comparações e contrastes em diferentes situações ”, “ Detectar contradições ”, “ Fazer escolhas valorativas avaliando consequências ”, “ Questionar a realidade “ e “ Argumentar coerentemente ”.
O Gráfico 3.7 mostra, a título de ilustração, o comportamento das respostas à questão 2 de Formação Geral para os estudantes concluintes da área de CIÊNCIAS ECONÔMICAS como função da nota geral na prova.
Trata -se de uma questão considerada fácil para os concluintes de CIÊNCIAS ECONÔMICAS e de dificuldade média para o conjunto de concluintes de todas as áreas.
O poder de discriminação foi médio para os estudantes da área e médio para o conjunto de todos os concluintes.
Neste Gráfico, cada uma das cinco curvas, representou o percentual de respostas em determinada alternativa da questão, em função do escore total dos estudantes concluintes da área de CIÊNCIAS ECONÔMICAS.
A curva vermelha corresponde à alternativa C, a correta para esta questão.
Assim, observa -se que entre os estudantes com notas mais baixas – desempenho geral mais fraco na prova como um todo –, alem da escolha da alternativa correta, houve a escolha de alternativas incorretas, especialmente a alternativa E – curva \textbf{roxa}, seguida da alternativa B – curva em verde claro.
A proporção de estudantes que selecionaram a resposta correta C, aumentou com o desempenho chegando a 100 \%.
Os Gráficos relativos às demais questões de Formação Geral constam do Anexo I. 3.2.2 Componente Específico A Tabela 3.7 apresenta as estatísticas básicas em relação às questões objetivas do Componente Específico.
Como pode ser observado, os estudantes concluintes obtiveram desempenho superior, com uma média de 39,40 pontos, em relação aos estudantes ingressantes com uma média de 30,66 pontos.
Os desvios-padrão indicaram que a variabilidade entre os concluintes foi maior ( 18,6 ) do que entre os ingressantes ( 15,3 ).
Os coeficientes de assimetria para os dois grupos foram positivos indicando uma inclinação para à esquerda.
A assimetria foi, porém, maior entre os ingressantes.
A análise das questões objetivas do Componente Específico em CIÊNCIAS ECONÔMICAS, em relação ao grau de facilidade da prova ( ver Tabela 3.8 ), mostrou que, das 27 questões, 20 encontraram -se na categoria difícil ( entre 15 \% e 40 \%, não incluso, de acertos ), 6 foram de nível médio ( entre 40 \% e 60 \%, não incluso, de acertos ), nenhuma está na categoria fácil ( entre 60 \% e 85 \%, não incluso, de acertos ) e na classificação muito fácil ( a partir de 85 \% de acertos ).
Uma questão teve percentual de acertos menor do que 15 \% ( classificação muito difícil ).
Nenhuma questão foi desconsiderada pela Comissão Assessora da Área de CIÊNCIAS ECONÔMICAS.
Das 27 questões objetivas relativas ao Componente Específico ( ver Tabela 3.9 ), 6 obtiveram o índice de discriminação muito bom, 6 obtiveram o índice de discriminação bom e 3 obtiveram nível de discriminação médio.
O índice fraco que eliminaria a questão, foi alcançado por 12 questões.
Na sequência, no Quadro 3.2 estão especificadas, por questão, as habilidades aferidas pelas questões do Componente Específico da prova, bem como os resultados da análise relativa aos índices de facilidade e de discriminação de cada questão.
O Gráfico 3.8 mostra, a título de ilustração, o comportamento da questão 34 do Componente Específico para os estudantes concluintes da área de CIÊNCIAS ECONÔMICAS como função da nota geral na prova5.
Trata -se de uma questão considerada difícil para os concluintes de CIÊNCIAS ECONÔMICAS.
O poder de discriminação foi muito bom para os estudantes da área.
Neste Gráfico, cada uma das cinco curvas, representa o percentual de respostas em determinada alternativa da questão, em função da nota total dos estudantes concluintes de CIÊNCIAS ECONÔMICAS.
A curva verde escuro correspondeu à alternativa D, a correta para esta questão.
Assim, observou -se que entre os estudantes com escores mais baixos – desempenho geral mais fraco na prova como um todo –, houve a escolha de uma alternativa incorreta, a alternativa C – curva vermelha.
A proporção de estudantes que selecionaram a resposta correta C aumentou com o desempenho deste grupo de estudantes chegando ao acerto de 100 \% dos estudantes com notas pelo menos 87,5.
Esta análise permitiu verificar como a questão discriminou entre os grupos de desempenho.
Os Gráficos relativos às demais questões do Componente Específico constam do Anexo I. 3.3 Análise das Questões Discursivas 3.3.1 Formação Geral A Tabela 3.10 apresenta as estatísticas básicas da prova por grupo de estudantes.
A análise dos resultados de desempenho dos estudantes nas questões discursivas que são apresentados na tabela 3.10 e no gráfico 3.9, mostrou que as notas foram mais baixas no conjunto dessas questões que no das questões objetivas.
O total de estudantes inscritos para a prova foi de 20821.
Destes, 75,3 \% compareceram, sendo que a presença foi maior entre concluintes ( 83,1 \% ) do que entre ingressantes ( 69,4 \% ).
A média geral da prova foi de 27,50, sendo que os ingressantes obtiveram média mais baixa, de 26,14, que os concluintes, cuja média foi 29,02.
O desvio-padrão geral foi de 34,9, sendo o do grupo de ingressantes menor ( 34,4 ), em relação ao grupo dos concluintes ( 35,4 ), aqueles possuíram uma distribuição mais homogênea nas notas.
A nota máxima foi 100,0, atingida por concluintes e ingressantes.
Os coeficientes de assimetria para os dois grupos considerados, o de concluintes e o de ingressantes, bem como para a população como um todo foram todos positivos, indicando uma maior concentração à esquerda da distribuição.
Para todos os grupos pouco mais da metade dos estudantes presentes obtiveram nota zero nas questões discursivas da Formação Geral.
Das 8.915 provas com nota zero, 360 ( 4,0 \% ) continham um protesto, 2596 ( 29,1 \% ) estavam em branco e 5959 ( 66,8 \% ) foram respondidas, porém de forma inadequada.
Entre os concluintes a maioria ( 70,5 \% ) dentre aqueles com nota zero apresentou a prova em branco.
Entre os ingressantes 63,8 \% das notas nulas ( zero ) decorreram do desempenho do próprio estudante verificadas na correção.
Uma comparação entre o desempenho de ingressantes e concluintes pode ser realizada a partir dos dados apresentados no Gráfico 3.9.
As notas dos concluintes, de um modo geral, foram ligeiramente mais elevadas do que as dos ingressantes, cuja distribuição de notas está à direita dos primeiros.
A ideia de melhor desempenho foi reforçada pela proporção de notas nulas : 54,9 \% entre os concluintes e 59,0 \% entre os ingressantes.
Na sequência, discorrer -se-á, especificamente, sobre uma das questões discursivas de Formação Geral, a número 10 já que a número 9 foi desconsiderada.
Em especial, serão analisados os conteúdos englobados na questão e os resultados obtidos pelos estudantes.
Ao final, são apresentados os comentários da banca corretora a respeito do observado nas respostas dos estudantes, suas impressões e conclusões.
Cumpre esclarecer que, tendo em vista que as questões discursivas de Formação Geral foram padronizadas, ou seja, constam de todas as provas, os comentários da banca foram os mesmos para todas as áreas acadêmicas, eis que direcionados a todos os estudantes que participaram do ENADE/2009. 3.3.1.1 Análise da Questão Discursiva 9 de Formação Geral DESCONSIDERADA 3.3.1.2 Análise da Questão Discursiva 10 de Formação Geral A Questão Discursiva nº. 10 buscou avaliar as habilidades : • Ler e interpretar textos ; • Analisar e criticar informações ; • Extrair conclusões por indução e/ou dedução ; • Estabelecer relações, comparações e contrastes em diferentes situações ; • Fazer escolhas valorativas avaliando consequências ; • Questionar a realidade ; • Argumentar coerentemente.
O Quadro 3.3 apresenta as habilidades/competências, temas predominantes e estatísticas básicas da questão 10 de Formação Geral por grupo de estudantes.
Cumpre notar que como a questão 9 foi desconsiderada, a nota referente às questões discursivas da Formação Geral é idêntica à nota da questão 10.
A média geral desta questão foi de 27,50, sendo que os ingressantes obtiveram média mais baixa, de 26,14, que os concluintes, cuja média foi 29,02.
O desvio-padrão geral foi de 34,9, sendo que o do grupo de ingressantes foi ligeiramente menor ( 34,4 ), em relação ao grupo dos concluintes ( 35,4 ), aqueles possuíram uma distribuição mais homogênea nas notas.
A nota máxima foi 100,0, atingida por concluintes e ingressantes.
Os coeficientes de assimetria para os dois grupos considerados, o de concluintes e o de ingressantes, bem como para a população como um todo foram positivos, indicando uma maior concentração à esquerda da distribuição.
Para todos os grupos, cerca de metade dos estudantes presentes obtiveram nota zero nesta questão discursiva de Formação Geral.
Das 8.915 provas com nota zero, 360 ( 4,0 \% ) continham um protesto, 2596 ( 29,1 \% ) estavam em branco e 5959 ( 66,8 \% ) foram respondidas, porém de forma inadequada.
Entre os concluintes a maioria ( 70,5 \% ) dentre aqueles com nota zero estava com a prova em branco.
Entre os ingressantes, 63,8 \% das notas nulas ( zero ) decorreram do desempenho do próprio estudante verificadas na correção.
Uma comparação entre o desempenho de ingressantes e concluintes para esta questão foi exatamente igual à já realizada para a nota discursiva de Formação Geral a partir dos dados apresentados no Gráfico 3.9. 3.3.1.3 Comentários sobre as respostas à Questão Discursiva 10 Apresentação A questão 10 de Formação Geral apresentou um questionamento capaz de analisar a habilidade do estudante em reconhecer o que é um direito social, optar por um direito social vigente que possa diminuir as desigualdades de renda familiar no Brasil e discutir os argumentos em favor do direito social escolhido como suporte à resposta.
Para análise das respostas apresentadas pelos estudantes, houve uma leitura inicial bastante criteriosa.
Foram considerados para efeito de pontuação, os seguintes argumentos : • Acesso à educação pública, gratuita e de qualidade, garantindo ao cidadão melhores condições de vida ; • Assistência à saúde, em seu contexto mais amplo, garantindo os cuidados essenciais de saúde, baseados em técnicas apropriadas ao alcance de todas as pessoas, \textbf{contando} com a participação da população ; • Permanência do estudante na escola, em todos os níveis escolares ( da educação infantil à educação superior, possibilitando a qualificação profissional ), o que garante melhores condições de trabalho e remuneração.
Considerou -se a qualificação e o desempenho docente, desenvolvendo nos estudantes melhores expectativas de aprendizagem e interesse pelo conteúdo da escola.
Os professores desenvolvendo metodologias adequadas e conhecimento pleno quanto a questões de aprendizagem ; • Condições dignas de trabalho, com remuneração que garanta qualidade de vida.
Foram consideradas as condições de trabalho, a preparação profissional para o trabalho e os cursos de aperfeiçoamento, oportunizando momentos de aprendizagem e desenvolvimento, a partir de situações desencadeadoras ; • Ser proprietário do imóvel que reside, reduzindo os gastos com aluguel, garantindo melhores condições de vida, equilíbrio financeiro e construção do bem-estar familiar.
Desempenho dos estudantes na questão É importante salientar que havia respostas que desenvolveram adequadamente o questionamento apresentado, com fundamentação baseada em argumentos coerentes.
Contudo, em um número considerável de respostas, havia apenas um argumento e não dois, conforme fora solicitado.
Muitos estudantes informaram sobre problemas da educação, de um modo geral.
Algumas observações mencionadas pelos corretores durante o processo de correção foram : • há preocupação com uma educação de qualidade, com sugestões de cursos profissionalizantes e métodos preparatórios de aulas. • a educação é revelada como de baixa qualidade, exaltando -se treinamentos e cursos de atualização para a classe docente. • a saúde como essência para a vida.
O SUS é colocado como parte integrante de direito social, garantindo às pessoas fácil acesso. • oportunidade de emprego como forma de ascensão social, direito, cidadania. • os políticos brasileiros como causadores do despreparo educacional.
As respostas analisadas apresentaram norma culta, coesão e coerência razoáveis, que garantiram a compreensão da mensagem desenvolvida pelos estudantes.
Os erros ortográficos foram poucos e não comprometeram a discussão sobre o assunto e sobre os pontos de vista abordados.
Quanto à coerência, houve um desvio quanto ao objetivo da questão, já que alguns estudantes, através dos seus argumentos retrataram reivindicações, reclamações e situações de problemas enfrentados pelos brasileiros, citando golpes políticos.
Houve erros quanto à concordância verbal e nominal.
Houve, também, dificuldade em utilizar os conectores corretamente.
Alguns estudantes apresentaram em suas respostas redundância e prolixidade.
Houve erros quanto à concordância verbal e nominal.
Houve também, dificuldade em utilizar os conectores corretamente.
O desempenho foi considerado satisfatório, pois a maioria dos estudantes soube responder ao questionamento.
Houve domínio do conteúdo, pois as ideias, os argumentos e a originalidade corresponderam ao esperado. 3.3.2 Componente Específico Nas questões discursivas do Componente Específico – questões de números 38, 39 e 40 –, a diferença de desempenho entre os grupos não foi muito acentuada, inclusive menor do que para as questões discursivas de Formação Geral.
A Tabela 3.11 contém as estatísticas da média destas questões.
Observa -se que no grupo dos concluintes a nota média alcançada foi 17,88 enquanto que no grupo dos ingressantes foi de 9,73, ambas menores do que as obtidas nas questões discursivas de Formação Geral, que foram de 26,14 para os ingressantes e 29,02 para os concluintes ( Tabela 3.10 ).
As médias foram também menores dos que as obtidas para as questões objetivas do Componente Específico ( Tabela 3.7 ) : respectivamente 30,66 e 39,40 para ingressantes e concluintes.
A análise de cada uma destas questões será feita nas subseções seguintes.
Os valores positivos para a medida de assimetria indicaram uma distribuição com cauda para a direita.
Dos 15.681 presentes, 54,1 \% obtiveram nota zero.
A tabela 3.12 apresenta um detalhamento da situação destes estudantes.
Do total de pessoas que compareceram ( 41,6 \% ) não fez nenhuma das 3 questões.
Esta proporção foi menor entre os concluintes 35,0 \% do que entre os ingressantes ( 46,2 \% ).
A combinação de questões em branco e preenchidas correspondeu a 97,4 \% das provas.
Menos de 1 \% no total das provas correspondeu a 3 questões com protesto.
O Gráfico 3.10 apresenta a distribuição das notas nas questões discursivas do Componente Específico por grupo de estudantes.
A moda se situou no primeiro intervalo, aquele com valores entre 0 e 10 incluindo os extremos do intervalo.
A impressão de assimetria informada pelos coeficientes de assimetria positivos ( Tabela 3.11 ) foi reforçada no gráfico pela concentração no primeiro intervalo. 3.3.2.1 Análise da Questão Discursiva 38 do Componente Específico A questão 38, cujas habilidades aferidas encontram -se descritas no Quadro 3.4, que ainda contém outras informações, seguido da Tabela 3.13 e do Gráfico 3.11, foi aquela em que os grupos de estudantes observados obtiveram um desempenho pior entre as questões discursivas do Componente Específico.
A nota média do grupo dos ingressantes foi de 7,44, e a dos concluintes 17,98.
O percentual de notas zero obtidas nos dois grupos foi de 68,3 \%, maior entre os ingressantes ( 77,7 \% ) do que entre os concluintes ( 57,8 \% ).
Um pouco mais da metade dos ingressantes que tiraram zero apresentaram a prova estava em branco.
Entre os concluintes, quase 60 \% responderam a questão. ( ver Tabela 3.13 ).
Uma pequena minoria ( 1,7 \% ) utilizou a questão para protestar.
Apesar do baixo desempenho, em ambos os grupos a nota máxima de 100,0 pontos foi alcançada.
O Gráfico 3.11 apresenta a distribuição das notas na questão 38 discursiva do Componente Específico por grupo de estudantes.
A moda se situou no primeiro intervalo, aquele com valores entre 0 e 10 incluindo os extremos do intervalo.
A impressão de assimetria informada pelos coeficientes de assimetria positivos ( Quadro 3.4 ) foi reforçada no gráfico pela concentração no primeiro intervalo. 3.3.2.2 Comentários sobre as respostas à Questão Discursiva 38 A questão foi elaborada de forma concisa e clara quanto aos conceitos solicitados e às periodicidades dos fatos econômicos sobre os quais os conhecimentos seriam testados.
Isto atendeu plenamente às habilidades e competências requeridas, justapondo dois períodos : Plano Cruzado ( 1986 ) e Plano Real ( 1994 ).
A questão não apresentou dificuldades para estudantes dos últimos semestres, muito embora tenham apresentado dificuldades no contexto teórico e histórico dos fatos econômicos questionados.
Como análise do desempenho dos estudantes, temos a considerar que houve uma grande diferença na pontuação, sendo as mais baixas, dos ingressantes o que foi constatado após o final da correção, com conhecimento dos dados estatísticos.
Também há que considerar que a maioria dos estudantes, neste caso, concluintes, destacou pelo menos um dos principais conceitos referentes ao Plano Cruzado e do Plano Real e alguns outros chegaram a mencionar outros aspectos dos referidos planos, demonstrando conhecimento histórico da economia brasileira.
Assim sendo, percebeu -se que a questão foi bem elaborada e que os estudantes que a responderam fizeram pelo entendimento do enunciado. 3.3.2.3 Análise da Questão Discursiva 39 do Componente Específico A questão 39, cujas habilidades aferidas encontram -se descritas no Quadro 3.5, que ainda contém outras informações, seguido da Tabela 3.14 e do Gráfico 3.12, foi aquela em que os grupos de estudantes observados obtiveram o pior desempenho entre as questões discursivas do Componente Específico.
A nota média do grupo dos ingressantes foi de 8,84, e a dos concluintes 16,76, menor entre os ingressantes.
Isto ocorre, principalmente pelo grande percentual de notas nulas : 75,1 \% entre os ingressantes e 60,4 \% entre os concluintes.
Das provas com notas nulas nesta questão, uma grande proporção foi deixada em branco : 53,2 \% entre os ingressantes e 43,7 \% entre os concluintes ( ver Tabela 3.14 ).
Entre os com nota nula, uma pequena minoria ( 2,6 \% ) utilizou a questão para protestar.
Apesar do baixo desempenho, em ambos os grupos a nota máxima de 100,0 pontos foi alcançada.
Os coeficientes de assimetria positivos indicaram uma concentração para a esquerda e espalhamento a direita.
O Gráfico 3.12 apresenta a distribuição das notas na questão 39 discursiva do Componente Específico por grupo de estudantes.
A moda se situou no primeiro intervalo, aquele com valores entre 0 e 10 incluindo os extremos do intervalo.
A impressão de assimetria informada pelos coeficientes de assimetria positivos ( Quadro 3.5 ) foi reforçada no gráfico pela concentração no primeiro intervalo. 3.3.2.4 Comentários sobre as respostas à Questão Discursiva 39 A questão aborda um fato da atualidade e com padrão de resposta preciso, com grau de dificuldade mediano.
Para os estudantes concluintes, a questão não apresenta dificuldades, mas para os iniciantes é difícil, porque exige maior conhecimento em teoria macroeconômica.
Considerando a totalidade das correções, observou -se que houve um pequeno número de estudantes com brilhante desempenho, demonstrando conhecimento e maturidade no desenvolvimento da matéria apresentada.
Verificou -se também, que muitos estudantes relataram que a renda iria aumentar, sem no entanto explicitar que este aumento seria mais do que proporcional, fator caracterizado no Modelo Keynesiano Simples.
Conclui -se que as pontuações com disparidades ocorreram porque os estudantes ingressantes e concluintes não possuem o mesmo nível de formação. 3.3.2.5 Análise da Questão Discursiva 40 do Componente Específico A questão 40, cujas habilidades aferidas encontram -se descritas no Quadro 3.6, que ainda contém outras informações, seguido da Tabela 3.15 e do Gráfico 3.13.
A questão foi aquela em que os grupos de estudantes observados obtiveram o melhor desempenho entre as questões discursivas do Componente Específico.
A nota média do grupo dos ingressantes foi de 12,94, e a dos concluintes 18,95.
Nesta questão ocorre também o menor percentual de notas nulas : 70,0 \% entre os ingressantes e 59,6 \% entre os concluintes.
A maior parte das provas com nota zero decorreu de provas preenchidas : 47,1 \% entre os ingressantes e 54,5 \% entre os concluintes ( ver Tabela 3.15 ).
Uma pequena minoria ( 1,7 \% ), dentre aqueles com nota nula, utilizou a questão para protestar.
Em ambos os grupos a nota máxima de 100,0 pontos foi alcançada.
Os coeficientes de assimetria positivos indicaram uma concentração à esquerda.
O Gráfico 3.13 apresenta a distribuição das notas na questão 40 discursiva do Componente Específico por grupo de estudantes.
A moda se situou no primeiro intervalo, aquele com valores entre 0 e 10 incluindo os extremos do intervalo.
A impressão de assimetria informada pelos coeficientes de assimetria positivos ( Quadro 3.7 ) foi reforçada no gráfico pela concentração no primeiro intervalo. 3.3.2.6 Comentários sobre as respostas à Questão Discursiva 40 A questão em análise apresentou um enunciado baseado no conteúdo predominante da Formação Teórico-Quantitativa.
Trata -se de uma análise sobre os preços e quantidades com utilização de gráficos de demanda e oferta.
A questão atende aos conteúdos tratados no primeiro ano do curso, portanto com boas condições de acerto para estudantes ingressantes e concluintes, e ainda, oferece margem à interpretação sobre o aspecto de elasticidade dos produtos, alvo da análise.
Considerando tais aspectos, esta foi a questão com melhor pontuação pelos estudantes, devendo enfatizar que os que não a acertaram no todo, pelo menos apresentaram alguma tentativa sobre o problema proposto.
Houve, portanto, maior homogeneidade de notas, dos estudantes, nesta questão.
De modo geral, o padrão de resposta proposto foi observado pelos estudantes, o que determinou o melhor rendimento.

% matched lemmas: contar, roxa
\end{textsample}
